\documentclass[11pt]{article}
\usepackage{fontspec}
\usepackage[margin=1in,top=1.1in,bottom=1.1in]{geometry}
\usepackage{booktabs}
\usepackage{longtable}
\usepackage{array}
\usepackage{amsmath}
\usepackage{xcolor}
\usepackage{titlesec}
\usepackage{fancyhdr}
\usepackage{enumitem}
\usepackage{caption}
\usepackage{unicode-math}
\usepackage[hidelinks]{hyperref}
\usepackage{microtype}
\usepackage{tabularx}
\usepackage{colortbl}
\usepackage{needspace}
\usepackage{float}

\setmainfont{Times New Roman}
\setmonofont{Consolas}[Scale=0.86]
\setmathfont{Cambria Math}

\definecolor{dqblue}{HTML}{1F3864}
\definecolor{dqgrey}{HTML}{5A5A5A}
\definecolor{dqrule}{HTML}{B4B4B4}
\definecolor{rowtint}{HTML}{F2F2F2}
\definecolor{warnbg}{HTML}{FBF3E4}
\definecolor{winbg}{HTML}{EAF1E4}
\definecolor{ctlbg}{HTML}{F7E9E9}

\titleformat{\section}{\Large\bfseries\color{dqblue}}{\thesection}{0.7em}{}
\titleformat{\subsection}{\large\bfseries\color{dqblue}}{\thesubsection}{0.6em}{}
\titleformat{\subsubsection}{\normalsize\bfseries}{\thesubsubsection}{0.5em}{}
\titlespacing{\section}{0pt}{1.6ex plus .3ex}{0.8ex}
\titlespacing{\subsection}{0pt}{1.3ex plus .3ex}{0.6ex}

\setlength{\headheight}{21.3pt}
\addtolength{\topmargin}{-7.3pt}
\setlength{\emergencystretch}{3em}
\pagestyle{fancy}
\fancyhf{}
\fancyhead[L]{\footnotesize\color{dqgrey}DynQuant Phase 2 --- Beating GPTQ at 3 Bits}
\fancyhead[R]{\footnotesize\color{dqgrey}\thepage}
\renewcommand{\headrulewidth}{0.4pt}
\renewcommand{\headrule}{\hbox to\headwidth{\color{dqrule}\leaders\hrule height \headrulewidth\hfill}}

\captionsetup{font=small,labelfont={bf,color=dqblue}}
\setlist{itemsep=0.25ex,topsep=0.5ex,parsep=0pt}

\newcommand{\x}{\texorpdfstring{$\times$}{x}}
\newcommand{\pt}{\,pts}
\newcommand{\pct}{\,\%}
\newcommand{\code}[1]{\texttt{#1}}
\newcommand{\minus}{$-$}
\newcommand{\e}[2]{$#1\!\times\!10^{#2}$}
\newcommand{\best}[1]{\textbf{#1}}

\newsavebox{\asidebox}
\newenvironment{aside}%
{\par\needspace{4\baselineskip}\vspace{0.8ex}%
\begin{lrbox}{\asidebox}\begin{minipage}{\dimexpr\textwidth-18pt\relax}\small}%
{\end{minipage}\end{lrbox}%
\noindent\setlength{\fboxsep}{9pt}\colorbox{warnbg}{\usebox{\asidebox}}%
\vspace{0.6ex}\par}

\begin{document}

%% ============================================================ TITLE
\begin{titlepage}
\centering
\vspace*{1.0in}
{\color{dqblue}\rule{\textwidth}{1.2pt}}\\[0.8em]
{\Huge\bfseries DynQuant --- Phase 2\par}
\vspace{0.5em}
{\LARGE Beating GPTQ at 3 Bits\par}
\vspace{0.35em}
{\large Four stacked levers, no error feedback, no calibration set\par}
\vspace{0.7em}
{\color{dqblue}\rule{\textwidth}{1.2pt}}\\[1.8em]
{\large 17 experimental arms, 2 negative controls, 6 frontier rungs\par}
\vspace{0.3em}
{\large Qwen3.5-2B-Base / CaseHOLD, $n = 5{,}314$\par}
\vspace{2.0em}

\begin{minipage}{0.86\textwidth}
\small
\textbf{Summary.} Phase 1 closed with one unfavourable finding that replicated on every panel it
was measured on: \emph{GPTQ wins the 3-bit tier}. At matched bytes GPTQ was simultaneously
3.1\pct{} smaller and 1.34 points better. Phase 2 set out to reverse that on the panel where it
was measured most sharply, under a hard constraint --- \textbf{do not copy GPTQ's mechanism}. No
error feedback, no inverse Hessian, no sequential column compensation, and above all no
calibration set, since a fine-tune-time signal hook exists precisely to avoid needing one.
\\[0.7em]
It is reversed. \code{p2\_rb\_agg} scores \textbf{89.57\pct{} at 708{,}087{,}808 bytes} against
\code{gptq\_3b\_head}'s \textbf{88.03\pct{} at 741{,}475{,}927} --- \textbf{$+1.54$ points,
203/121 discordant, exact McNemar $p < 0.0001$, CI $[+0.88, +2.21]$, at 4.5\pct{} fewer bytes}.
It stays ahead at 8.3\pct{} fewer bytes ($+1.05$, $p = 0.0026$) and level at 16.4\pct{} fewer
($-0.19$, $p = 0.64$). The final arm gives up \textbf{0.17 points against full precision at 3.01
stored bits}.
\\[0.7em]
\textbf{Settled (2026-08-15), by a replicate.} The \code{gptq\_3b\_head} arm above was fitted on
a \emph{symmetric} grid and every DynQuant arm on an asymmetric one. That was not a choice this
report made and recorded; it was hardcoded in \code{experiments/\_llmc.py} from that file's first
commit --- the recipe builder asked for a symmetric grid for every method except AWQ --- so no
GPTQ arm this project had published had ever been fitted the other way. On the phase-4
Mistral-7B panel the same difference, holding method, calibration set, group size and byte
anchor fixed, was worth \textbf{$+69.4$ points}. This report held its own headline pending the
matching control on \emph{this} panel. The original checkpoint no longer exists, so the control
was run as a \textbf{replicate}: same recipe, a fresh fine-tune, torch 2.12.0+cu130 /
transformers 5.10.1 in place of the original's torch 2.13 / transformers 5.14. The real
asymmetric arm scored \textbf{88.45\pct{} at 741{,}486{,}130 bytes}, and DynQuant still won:
\textbf{$+1.13$, 155/95 discordant, exact $p = 0.000178$, CI $[+0.55, +1.71]$, at 4.5\pct{} fewer
bytes.} Against the replicate's symmetric arm (87.67\pct{} at 735{,}981{,}792 bytes) the margin is
$+1.90$ ($p = 2.0\mathrm{e}{-8}$) at 3.79\pct{} fewer bytes, and against the replicate's own bf16
ceiling (89.42\pct{}) DynQuant is a statistical tie ($+0.15$, $p = 0.57$) at $5.32\times$ fewer
bytes. Full account in \S\ref{sec:addendum}. Absolute figures in this replicate differ from the
panel above at the decimal --- different checkpoint, different library versions --- so the two
are not merged; every comparison inside \S\ref{sec:addendum} is arm-vs-arm within the replicate's
own run, as everywhere else in this report.

The margin comes from four levers, each isolated at byte-identical budgets: a Gauss--Newton
sensitivity model, a row-partitioned tied embedding, an $\mathbb{E}[x^2]$-weighted clip
objective, and per-row body allocation. The last of those is also the report's central finding,
and it is not the flattering reading of itself: \textbf{allocation granularity is a multiplier on
the signal, not a gain of its own.} With the row widths shuffled, per-row allocation \emph{loses}
1.28 points at identical bytes; with the signal's ordering it wins 2.31 against that control.
\end{minipage}

\vfill
{\small\color{dqgrey}
Measured 2026-07-30 on a single NVIDIA A100 80\,GB. All arms score the same 5{,}314 items in the
same order, so every comparison is a paired exact McNemar test on stored per-item hits.\\
\textbf{Scope: one panel.} \S\ref{sec:notoverturn} states what this does \emph{not} overturn.\par}
\vspace{0.6em}
\end{titlepage}

\tableofcontents
\newpage

%% ============================================================ 1
\section{The result}

\subsection{What had to be reversed}

Phase 1's scorecard had one unambiguous row, and it was the unfavourable one. At a matched 3-bit
budget on Qwen3.5-2B-Base / CaseHOLD:

\begin{center}\small
\begin{tabular}{@{}lrrl@{}}
\toprule
Phase 1, matched bytes & delta & $p$ & \\
\midrule
\code{dq\_3p25} vs \code{gptq\_3b\_head} & \best{$-1.34$} & \best{\e{4.2}{-4}} & GPTQ 3.1\pct{} smaller \emph{and} 1.34 better \\
\bottomrule
\end{tabular}
\end{center}

It replicated on all three panels ($-2.28$ / $-1.04$ / $-0.78$ against the shipped-convention
GPTQ arms) and was described in the Phase 1 report as the most robust conclusion the campaign
produced. The diagnosis was equally clear: \emph{DynQuant chooses widths well and then rounds
naively; GPTQ chooses nothing and then rounds cleverly. At 3 bits the second matters more.}

\subsection{The headline}

\begin{center}\small
\begin{tabular}{@{}llrrrrr@{}}
\toprule
 & arm & acc \% & bytes & avg bits & vs \code{gptq\_3b\_head} & $p$ \\
\midrule
before & \code{dq\_3p25} (shipped) & 86.70 & 764{,}290{,}013 & 3.2494 & \best{$-1.34$} & \best{0.0004} \\
tie & \code{p2\_wc\_agg} (per-module body) & 88.54 & 708{,}087{,}808 & 3.0104 & $+0.51$ & 0.16 (ns) \\
\rowcolor{winbg}
\textbf{after} & \code{p2\_rb\_agg} (per-row body) & \best{89.57} & \best{708{,}087{,}808} & 3.0104 & \best{$+1.54$} & \best{$<0.0001$} \\
\bottomrule
\end{tabular}
\end{center}

\code{rb\_agg} vs \code{dq\_3p25} is $+2.87$ points end to end, at \textbf{7.4\pct{} fewer bytes}
than the shipped arm. The fp16 ceiling on this panel is 89.74\pct, so the final arm gives up
\textbf{0.17 points against full precision at 3.01 stored bits}.

Both endpoints are significant in their own right: GPTQ really did win before ($-1.34$,
$p = 0.0004$), and it really is beaten now ($+1.54$, 203/121 discordant, CI $[+0.88, +2.21]$).

\begin{aside}
\textbf{The middle row stays in the table.} At the time it was the best arm available and it was
$+0.51$ at $p = 0.16$ --- a tie --- and it was written up as a tie for several days while it was
the best result in hand. It never did become significant. The win came from a \emph{different
mechanism} landing underneath it, not from re-reading the same number more favourably. Had
$+0.51$ been promoted to ``beats'' then, the actual result would have had nothing left to report.
\end{aside}

\subsection{The constraint the solution had to satisfy}

The obvious fix for ``rounds naively'' is GPTQ's own fix: propagate the quantization error of each
column into the not-yet-quantized columns using an inverse Hessian estimated from calibration
activations. It was considered and rejected, for two reasons that are worth stating because they
shaped everything that follows.

\begin{enumerate}[leftmargin=1.4em]
\item \textbf{It would make the method a reimplementation.} Signal-driven allocation plus
GPTQ's compensation is GPTQ with a preprocessing step. The interesting question --- how much is
left on the table by DynQuant's \emph{own} components --- would go unanswered.
\item \textbf{It requires a calibration dataset.} Hessian estimation needs activation statistics
over a held-out corpus. The entire premise of the fine-tune-time hook is that those statistics are
already being produced for free by the training run, and that no separate calibration pass is
needed. Adopting error feedback would have reintroduced the dependency the method exists to
remove.
\end{enumerate}

So the brief for Phase 2 was: \textbf{fix ours}. Every lever below acts on a component DynQuant
already owns --- the sensitivity model, the allocation partition, or the encoder's clip objective
--- and every input to every lever is a moment the \code{transformers} fine-tune hook already
accumulates online.

%% ============================================================ 2
\section{Statistical method}

Every number in this report is an exact paired McNemar test on the stored per-item \code{hits}
arrays. All arms score the same 5{,}314 items in the same order, so the arms are paired.

This matters more than it sounds. The unpaired binomial standard deviation used earlier in the
campaign is $\sigma = 0.419$ points, so an unpaired comparison of two arms carries
$\sqrt{2}\,\sigma \approx 0.59$ points of noise. The paired standard error of a \emph{difference}
is
\[
\mathrm{SE} \;=\; \frac{\sqrt{\,b + c - (b-c)^2/N\,}}{N},
\]
where $b$ and $c$ are the discordant counts, and here it runs 0.21--0.37 points --- roughly half
the unpaired figure, because two arms built from the same checkpoint agree on the overwhelming
majority of items and the pairing removes exactly that shared variance. The reported $p$ is a
two-sided exact binomial test on $b$ against $c$; the CI is the normal interval on the paired
difference.

\begin{aside}
Nothing was re-run to get this. The sharper test came free from per-item hit arrays already
written to disk. Its practical consequence cuts both ways: it promoted a $+1.07$ lever to
$p = 0.0002$, and it \emph{refused} to promote a $+0.60$ comparison against GPTQ's larger
configuration, which lands at $p = 0.069$ and is reported as a tie (\S\ref{sec:discipline}).
\end{aside}

%% ============================================================ 3
\section{The four levers}

Each lever was isolated at a \textbf{byte-identical} budget, so allocation or encoding is the only
difference between the two sides of every A/B. None of them is GPTQ's mechanism.

\begin{center}\footnotesize
\begin{tabular}{@{}llrrr@{}}
\toprule
lever & A/B & $\Delta$ & $b/c$ & $p$ \\
\midrule
Row-partitioned tied embedding vs uniform 2b & \code{tie\_u2\_ctl} $\to$ \code{rows\_agg}, $+7.00$ MiB & \best{$+2.90$} & 245/91 & \best{$<0.0001$} \\
Gauss--Newton sensitivity vs rank-product & \code{rows\_agg} $\to$ \code{combo\_agg} @ 708.1 MB & \best{$+1.07$} & 141/84 & \best{0.0002} \\
\textbf{Per-row body allocation vs per-module} & \code{wc\_agg} $\to$ \code{rb\_agg} @ 708.1 MB & \best{$+1.04$} & 122/67 & \best{0.0001} \\
$\mathbb{E}[x^2]$-weighted clip (encoder only) & \code{combo\_iso} $\to$ \code{wclip\_enc} @ 739.5 MB & $+0.70$ & 147/110 & 0.025 \\
\quad --- replication on a second map & \code{combo\_agg} $\to$ \code{wc\_agg} @ 708.1 MB & $+0.41$ & 135/113 & 0.18 \\
\quad --- \textbf{pooled over both replications} & & & 282/223 & \best{0.0098} \\
\midrule
\emph{(null)} re-pricing on weighted encoding & \code{wc\_agg} $\to$ \code{wc\_both\_agg} & $-0.09$ & 18/23 & 0.53 \\
\bottomrule
\end{tabular}
\end{center}

\subsection{Lever 1 --- Gauss--Newton sensitivity allocation}

The shipped scorer ranked each module on a product of two percentile ranks --- activation
saliency and gradient plasticity --- and fed the rank product to a greedy ROI knapsack. A rank is
ordinal: it says module $A$ matters more than module $B$, but not by how much, and not how much
worse $A$ gets between 3 bits and 2.

The replacement scores in units of loss. From the two per-channel second moments the hook already
accumulates --- $\mathbb{E}[x_c^2]$ over input channels and $\mathbb{E}[\delta_r^2]$ over
output-gradient channels --- the loss increase from quantizing module $m$ at width $b$ is the
diagonal Gauss--Newton (equivalently, empirical-Fisher) estimate
\begin{equation}
\mathcal{S}_m(b) \;=\; \sum_{r,c} \mathbb{E}[\delta_r^2]\,\mathbb{E}[x_c^2]\,
\bigl(W_{rc} - Q_b(W)_{rc}\bigr)^2 .
\label{eq:sens}
\end{equation}
The error term is not a surrogate curve: the weight block is run through the \emph{real} encoder
at the \emph{real} candidate width, clip search included. Allocation then becomes a
multiple-choice knapsack over $\{2,3,4,8\}$ in which the price of a move is the difference between
two sensitivities --- what that bit actually buys --- rather than a rank.

\textbf{Measured, byte-identical, one substitution:} \code{rows\_agg} 87.05\pct{} $\to$
\code{combo\_agg} 88.13\pct{} at the same 708{,}087{,}808 bytes. \best{$+1.07$ points, 141/84
discordant, $p = 0.0002$}. This is the cleanest single result in the phase: one component swapped,
byte-identical output.

\subsection{Lever 2 --- Row-partitioned tied embedding}

On Qwen3.5-2B the embedding and the LM head are the same $[248320, 2048]$ tensor, 27.05\pct{} of
the model. The shipped allocator gives that tensor one width. It does not have to: the format
already stores one fp16 scale and offset \emph{per row per group}, so splitting the rows into
width classes costs nothing extra in metadata.

The plan is \code{8b:2048, 4b:8192, 2b:*} --- the 2{,}048 most important rows by the signal at 8
bits, the next 8{,}192 at 4, everything else at 2.

\textbf{Measured against a flat 2-bit tie control} (\code{tie\_u2\_ctl} $\to$ \code{rows\_agg},
with the rank-product body held fixed on both sides so the comparison isolates the tie and
nothing else): \best{$+2.90$ points, 245/91, $p < 0.0001$}, for 7.00 MiB. That is 0.41 points per
MiB against the body's 0.011 --- roughly \textbf{37\x} the return per byte.

\textbf{This is the lever GPTQ structurally cannot pull}, and it is why the size gap exists at all.
See \S\ref{sec:tie}.

\subsection{Lever 3 --- $\mathbb{E}[x^2]$-weighted clip objective}

The encoder's per-group clip search minimised plain reconstruction error
$\lVert W - Q(W) \rVert^2$. But the quantity that matters is the error in the layer's
\emph{output}, and an input channel that carries large activations contributes proportionally
more to it. The objective becomes
\[
\min_{\alpha} \; \sum_{r,c} \mathbb{E}[x_c^2]\,\bigl(W_{rc} - Q_\alpha(W)_{rc}\bigr)^2 ,
\]
using the same per-channel second moment the sensitivity model already consumes. The clip grid was
also deepened, its floor moved from $\alpha = 0.80$ to $0.40$.

\textbf{Measured twice at byte parity, encoder only:} $+0.70$ ($p = 0.025$) on one map and
$+0.41$ ($p = 0.18$) on another. Neither replication clears 5\pct{} alone; \textbf{pooled over both,
282 vs 223 discordant, $p = 0.0098$.} The mechanism separates even though one instance of it does
not, and it is reported that way rather than by quoting the favourable replication.
\S\ref{sec:clip} decomposes where in the pipeline each clip change pays.

\subsection{Lever 4 --- Per-row body allocation}

The same idea as lever 2, applied to the body instead of the embedding: allocate a width to each
of the model's \textbf{571{,}968 body rows} rather than to each of its 187 modules. Same solver,
same sensitivity table, same 5{,}664{,}702{,}464-bit budget, same encoder, same clip grid, same
group size. Only the decision variable changes.

\textbf{Measured, byte-identical:} \code{wc\_agg} 88.54\pct{} $\to$ \code{rb\_agg} 89.57\pct{} at
708{,}087{,}808 bytes. \best{$+1.04$ points, 122/67, $p = 0.0001$}. This is the lever that turned
the tie into a win.

It also has the most misleading $\Delta$ in the table, and the next section is entirely about why.

%% ============================================================ 4
\section{Granularity is a multiplier on the signal, not a gain of its own}
\label{sec:granularity}

Reporting lever 4 as ``row allocation pays $+1.04$'' would have been wrong about which component
paid.

\subsection{The control}

The control permutes each module's width vector under a fixed seed. The width \emph{histogram} is
preserved exactly --- so the byte count is identical by construction, not by tuning --- and only
the question of \emph{which} rows get the wide ones changes. 196{,}071 of 571{,}968 rows
(34.3\pct) end up at a different width.

\begin{center}\small
\begin{tabular}{@{}lrrrr@{}}
\toprule
body allocation at 708{,}087{,}808 B & acc \% & vs module body & $b/c$ & $p$ \\
\midrule
per-module widths (\code{p2\_wc\_agg}) & 88.54 & --- & & \\
\rowcolor{ctlbg}
per-row, \textbf{shuffled within module} (\code{p2\_rb\_shuf}) & 87.26 & \best{$-1.28$} & 125/193 & \best{0.0002} \\
\rowcolor{winbg}
per-row, ordered by sensitivity (\code{p2\_rb\_agg}) & \best{89.57} & \best{$+1.04$} & 122/67 & \best{0.0001} \\
\bottomrule
\end{tabular}
\end{center}

\code{rb\_agg} vs \code{rb\_shuf} --- the signal's ordering, with everything else including the
histogram held fixed --- is \best{$+2.31$, 228/105, $p < 0.0001$, CI $[+1.64, +2.98]$}. The
decomposition is additive to 0.01 points: $-1.28$ (structure) $+$ 2.31 (signal) $= +1.03$ against
a measured $+1.04$.

\begin{aside}
\textbf{Row granularity on its own is negative.} Going from 187 decisions to 571{,}968 buys no
structural advantage --- it buys \emph{more decisions}, and more decisions made badly are worse
than fewer. This is the opposite of the natural intuition, which is that finer granularity must be
close to free because a module's rows are heterogeneous.
\end{aside}

They \emph{are} heterogeneous: the width the module allocator picks covers only \textbf{42.5\pct{}
of that module's rows at the median}, and under half the rows in \textbf{99 of 186 modules}. The
heterogeneity is real, and exploiting it still requires being right about which rows are which.

\subsection{What the row allocator does with the freedom: a barbell}

Same solver, same table, same 5{,}664{,}702{,}464-bit budget --- only the decision variable
changes:

\begin{center}\small
\begin{tabular}{@{}lrrr@{}}
\toprule
width & per-row params & per-module params & delta \\
\midrule
2 b & 628.0 M & 390.1 M & $+237.9$ M \\
3 b & 280.8 M & 612.4 M & \best{$-331.6$ M} \\
4 b & 424.6 M & 367.0 M & $+57.6$ M \\
8 b & 39.4 M & 3.3 M & \best{$+36.1$ M (12\x)} \\
\bottomrule
\end{tabular}
\end{center}

A per-module width is forced to the \emph{mean} of its rows' sensitivities; a per-row width serves
the \emph{distribution}. Loss is convex in width, so the optimum protects a small set of rows at 8
bits --- twelve times as many parameters as the module map could afford there --- and pays for it
by dumping the bulk from 3 to 2.

Per role: \code{self\_attn.k\_proj} gains 0.87 average bits, \code{o\_proj} 0.74,
\code{mlp.down\_proj} 0.51, while \code{linear\_attn.in\_proj\_qkv} gives up 0.59 and the tiny
\code{linear\_attn.in\_proj\_a} / \code{in\_proj\_b} give up 4.06 and 2.50 --- the module allocator
had those pinned at an 8-bit floor that most of their rows turn out not to need.

\subsection{Two costs the tables do not show}

\textbf{The per-row width code is not in the byte accounting.} The accounting charges payload plus
one fp16 scale and offset per row per group --- partition-invariant, so the A/Bs are honest --- but
not the width code itself: 4 distinct widths is 2 bits $\times$ 571{,}968 rows $=$
\textbf{142{,}992 B, 0.0202\pct{} of the checkpoint, 0.0006 average bits}. Three orders of
magnitude too small to explain any result here, but it is a real debit and it is stated rather than
buried.

\textbf{The arm was audited for materialization, and the audit was not optional.} 89.57\pct{} sits
0.17 under the fp16 ceiling and 0.13 \emph{above} the tie-only arm that leaves the body in bf16
(\code{p2\_tieonly}, 89.44\pct) --- which is exactly where a body that silently never got quantized
would land. \code{nbytes} is computed analytically from the plan, so the byte count cannot rule
that out. The arm was therefore re-encoded and inspected directly, counting distinct values inside
each group:

\begin{center}\small
\begin{tabular}{@{}lll@{}}
\toprule
row width & distinct values per group & relative reconstruction error \\
\midrule
2 b & exactly 4 & 0.0090 (an all-8b \code{out\_proj}) \\
3 b & exactly 8 & \quad \dots \\
4 b & exactly 16 & \quad \dots \\
8 b & 80--99 (of a possible 256) & 0.3799 (\code{in\_proj\_qkv}, layer 18) \\
\bottomrule
\end{tabular}
\end{center}

\textbf{Verdict: materialized.} No arm in this report is believed without this check.

\begin{aside}
\textbf{Operating rule: never ship a granularity change without the shuffled control.} It is one
extra evaluation, it preserves the byte count exactly by construction, and here it converted ``row
allocation is worth $+1.04$'' into ``the signal is worth $+2.31$ and the structure costs 1.28'' ---
a different claim about a different component. The same control is the honest test for any future
partitioning: per-group, per-column, per-expert.
\end{aside}

\subsection{What this does to Phase 1's headline attribution}

Phase 1 measured the training signal at \textbf{12\pct{} of the margin} over byte-matched RTN ---
22.62 points of allocator against 3.16 points of signal --- and concluded that the architectural
prior does most of the work. That measurement was taken at \emph{module} granularity, where there
are only 187 decisions and the architecture dominates them.

At row granularity the same signal is handed 571{,}968 decisions, and the shuffled control shows
its ordering carrying $+2.31$ of a $+1.04$ net. Both readings are true of what they measured, and
they are consistent: \textbf{granularity sets how much room the signal has to be right in.} Quote
the one that matches the granularity being shipped.

Note also what this rules out: the gain is \textbf{not} an artifact of finer \emph{encoding}. Every
arm here shares one encoder, one clip grid and one group size. Only the width vector differs.

%% ============================================================ 5
\section{The frontier}
\label{sec:frontier}

The recipe was walked down in roughly 30 MB steps to find where it breaks.
\code{gptq\_3b\_head} sits at 88.03\pct{} / 741{,}475{,}927 B throughout. Both descents are given,
because the second relocates the cliff rather than merely lifting the curve, and that is only
visible with the first one next to it.

\subsection{Per-module body --- the recipe as it stood before row allocation}

\begin{center}\small
\begin{tabular}{@{}lrrrrrrr@{}}
\toprule
arm & acc \% & bytes & avg bits & vs GPTQ & $p$ & vs rung above & $p$ \\
\midrule
\code{p2\_wc\_agg} & 88.54 & 708{,}087{,}808 ($-4.5\%$) & 3.0104 & $+0.51$ & 0.16 & --- & \\
\code{p2\_wc\_680} & 88.22 & 679{,}776{,}256 ($-8.3\%$) & 2.8900 & $+0.19$ & 0.63 & $-0.32$ & 0.14 \\
\code{p2\_wc\_650} & 87.22 & 649{,}891{,}840 ($-12.3\%$) & 2.7630 & \best{$-0.81$} & \best{0.034} & \best{$-1.00$} & \best{0.0004} \\
\code{p2\_wc\_620} & 86.41 & 619{,}876{,}352 ($-16.4\%$) & 2.6354 & \best{$-1.62$} & \best{$<0.0001$} & \best{$-0.81$} & \best{0.0058} \\
\bottomrule
\end{tabular}
\end{center}

The floor here is 680 MB. $708 \to 680$ is free ($-0.32$, only 68/51 discordant, not separated);
$680 \to 650$ is a cliff ($-1.00$, $p = 0.0004$), and 650 is the first rung that loses to GPTQ
significantly. The defensible statement was \emph{GPTQ's 3-bit accuracy at 8.3\pct{} fewer bytes},
with 12.3\pct{} measurably too far.

\subsection{Per-row body --- same targets, same pricer, same encoder}

\begin{center}\small
\begin{tabular}{@{}lrrrrrrr@{}}
\toprule
arm & acc \% & bytes & avg bits & vs GPTQ & $p$ & vs rung above & $p$ \\
\midrule
\rowcolor{winbg}
\code{p2\_rb\_agg} & \best{89.57} & 708{,}087{,}808 ($-4.5\%$) & 3.0104 & \best{$+1.54$} & \best{$<0.0001$} & --- & \\
\rowcolor{winbg}
\code{p2\_rb\_680} & \best{89.09} & 680{,}000{,}000 ($-8.3\%$) & 2.8910 & \best{$+1.05$} & \best{0.0026} & $-0.49$ & 0.014 \\
\code{p2\_rb\_650} & 88.61 & 649{,}999{,}872 ($-12.3\%$) & 2.7634 & $+0.58$ & 0.11 & $-0.47$ & 0.027 \\
\code{p2\_rb\_620} & 87.84 & 619{,}999{,}744 ($-16.4\%$) & 2.6358 & $-0.19$ & 0.64 & $-0.77$ & 0.0003 \\
\code{p2\_rb\_590} & 87.00 & 589{,}999{,}872 ($-20.4\%$) & 2.5084 & \best{$-1.04$} & \best{0.0065} & $-0.85$ & 0.0042 \\
\rowcolor{ctlbg}
\code{p2\_rb\_560} & 83.74 & 560{,}000{,}000 ($-24.5\%$) & 2.3809 & \best{$-4.29$} & \best{$<0.0001$} & \best{$-3.26$} & \best{$<0.0001$} \\
\bottomrule
\end{tabular}
\end{center}

\subsection{The gain grows as the budget tightens}

Row body against module body at byte parity, rung by rung:

\begin{center}\small
\begin{tabular}{@{}lrrrr@{}}
\toprule
budget & 708 MB & 680 MB & 650 MB $\dagger$ & 620 MB \\
\midrule
row $-$ module & $+1.04$ & $+0.87$ & $+1.39$ & $+1.43$ \\
$p$ & 0.0001 & 0.0022 & $<0.0001$ & 0.0001 \\
\bottomrule
\end{tabular}
\end{center}

$\dagger$ \textbf{The 650 rung is not a clean single-variable A/B.} \code{rb\_650} inherits its
clip metadata from the \code{wc\_650} map, which was priced with the weighted objective, but its
\emph{row} sensitivities were priced unweighted --- so that rung mixes granularity with pricing
objective. The 680 and 620 rungs are clean, and they bracket it at $+0.87$ and $+1.43$.

That ordering is the same mechanism as the barbell: the tighter the budget, the more the
module-uniform constraint costs, because forcing a whole module to the mean of its rows is exactly
the wrong move when there are not enough bits to go around.

\subsection{Three thresholds, all moved}

\begin{itemize}[leftmargin=1.4em]
\item \textbf{Beats GPTQ down to 680 MB} ($-8.3\%$), significantly, where the module recipe managed
$+0.19$ at $p = 0.63$.
\item \textbf{Ties GPTQ down to 620 MB} ($-16.4\%$, $-0.19$ at $p = 0.64$) --- a budget at which the
module recipe was losing by 1.62 with $p < 0.0001$.
\item \textbf{The cliff moved from 650 to below 590.} The module descent broke by a full point at
the first step past its floor; the row descent gives up 0.47--0.85 per 30 MB rung all the way down
to 590, then falls 3.26 in one step at 560.
\end{itemize}

So the claim the frontier supports is \textbf{GPTQ's 3-bit accuracy at 16.4\pct{} fewer bytes, and
better accuracy at 8.3\pct{} fewer}. The 560 MB rung is over the edge and it is reported because
\emph{a descent that never finds its own floor has not located anything}.

\subsection{Where the module cliff comes from, and an instrumentation gap}

The allocator's own diagnostics locate the module cliff: floors breached climb $60 \to 75 \to
\mathbf{92}$ of 187 modules across the first three rungs. The break arrives when about half the
model is pushed below its architectural floor --- which is also the point at which the soft-floor
path, the fix for the original allocator's silent \code{return floor\_map}, is doing nearly all the
work. That the curve degrades gracefully rather than collapsing is itself the soft-floor design
being vindicated: the shipped allocator would have returned the floor map unchanged and reported
nothing.

\begin{aside}
\textbf{The row descent has no equivalent diagnostic.} Role floors are defined per module and the
row allocator does not consult them. This is flagged as a gap rather than a feature: \textbf{the
reason 560 MB collapses is not currently instrumented.}
\end{aside}

%% ============================================================ 6
\section{Every arm}

The complete phase-2 ladder, all on Qwen3.5-2B-Base / CaseHOLD, $n = 5{,}314$, bf16 ceiling
89.74\pct.

\begin{center}\footnotesize
\begin{tabular}{@{}llrrr@{}}
\toprule
arm & what it is & acc \% & bytes & avg bits \\
\midrule
\rowcolor{winbg}
\code{p2\_rb\_agg} & final recipe: per-row body + row tie + weighted clip & \best{89.57} & 708{,}087{,}808 & 3.0104 \\
\code{p2\_rb\_680} & same recipe, 680 MB rung & 89.09 & 680{,}000{,}000 & 2.8910 \\
\code{p2\_rb\_650} & same recipe, 650 MB rung (mixed pricing $\dagger$) & 88.61 & 649{,}999{,}872 & 2.7634 \\
\code{p2\_wc\_agg} & per-module body, weighted clip encoder & 88.54 & 708{,}087{,}808 & 3.0104 \\
\code{p2\_wclip\_enc} & weighted clip, encoder only, iso budget & 88.45 & 739{,}545{,}088 & 3.1442 \\
\code{p2\_wc\_both\_agg} & weighted clip priced \emph{and} encoded & 88.45 & 708{,}087{,}808 & 3.0104 \\
\code{p2\_wc\_680} & per-module body, 680 MB rung & 88.22 & 679{,}776{,}256 & 2.8900 \\
\code{p2\_combo\_agg} & sensitivity allocator, unweighted clip & 88.13 & 708{,}087{,}808 & 3.0104 \\
\code{p2\_rb\_620} & per-row body, 620 MB rung & 87.84 & 619{,}999{,}744 & 2.6358 \\
\code{p2\_combo\_iso} & sensitivity allocator, iso budget & 87.75 & 739{,}545{,}088 & 3.1442 \\
\rowcolor{ctlbg}
\code{p2\_rb\_shuf} & \textbf{control} --- row widths shuffled within module & 87.26 & 708{,}087{,}808 & 3.0104 \\
\code{p2\_wc\_650} & per-module body, 650 MB rung & 87.22 & 649{,}891{,}840 & 2.7630 \\
\code{p2\_rows\_agg} & row tie + rank-product body & 87.05 & 708{,}087{,}808 & 3.0104 \\
\code{p2\_rb\_590} & per-row body, 590 MB rung & 87.00 & 589{,}999{,}872 & 2.5084 \\
\code{p2\_wc\_620} & per-module body, 620 MB rung & 86.41 & 619{,}876{,}352 & 2.6354 \\
\rowcolor{ctlbg}
\code{p2\_tie\_u2\_ctl} & \textbf{control} --- tie at uniform 2 b & 84.16 & 700{,}747{,}776 & 2.9792 \\
\code{p2\_rb\_560} & per-row body, 560 MB rung --- past the cliff & 83.74 & 560{,}000{,}000 & 2.3809 \\
\bottomrule
\end{tabular}
\end{center}

$\dagger$ See \S\ref{sec:frontier}: row sensitivities priced unweighted, clip metadata from a
weighted map.

\subsection{Reference arms}

\begin{center}\footnotesize
\begin{tabular}{@{}llrrr@{}}
\toprule
arm & what it is & acc \% & bytes & avg bits \\
\midrule
\code{bf16} & full precision ceiling & 89.74 & --- & 16 \\
\code{p2\_tieonly} & row-partitioned tie, body left in bf16 & 89.44 & --- & --- \\
\code{gptq\_3b} & llm-compressor, \code{ignore=["lm\_head"]}, tie left fp16 & 88.97 & 1{,}558{,}371{,}533 & 6.6253 \\
\code{gptq\_3b\_head} & llm-compressor, \code{ignore=[]} --- the like-for-like baseline & 88.03 & 741{,}475{,}927 & 3.1522 \\
\code{dq\_3p25} & DynQuant as shipped in Phase 1 & 86.70 & 764{,}290{,}013 & 3.2494 \\
\bottomrule
\end{tabular}
\end{center}

\code{p2\_tieonly} is the arm that makes the materialization audit in \S\ref{sec:granularity}
necessary: the final result sits 0.13 points \emph{above} it, which is precisely where an
unquantized body would land.

%% ============================================================ 7
\section{The clip objective pays in the encoder --- and the deep grid did not}
\label{sec:clip}

Both clip knobs were decomposed into (priced-old, encoded-new) against (priced-new, encoded-new).
They came out \textbf{opposite}, which is the useful part:

\begin{center}\small
\begin{tabular}{@{}lrr@{}}
\toprule
knob & encoder-only & $+$ re-pricing \\
\midrule
deep grid (floor $0.80 \to 0.40$) & $+0.09$ & \best{$+0.58$} \\
$\mathbb{E}[x^2]$-weighted objective & \best{$+0.41$} & $-0.09$ ($p = 0.53$) \\
\bottomrule
\end{tabular}
\end{center}

Re-pricing on top of the weighted encoder is \emph{nothing} --- 18 against 23 discordant items, the
smallest disagreement of any pair measured in the phase. Diffing the two maps says why:
\textbf{4 of 187 module widths changed.} The weighted objective cuts sensitivity 38.6\pct{} at 2 b
and 44.7\pct{} at 3 b, a near-common factor, and an allocator reads only \emph{ratios} --- between
modules and across widths --- so a common factor cancels and the ROI ordering survives.

The deep grid instead un-clamped only the low widths, moving \code{mlp.gate} L0's 2b/3b ratio from
4.69 to 3.47: roughly three times the ratio distortion, which is exactly where its $+0.58$ came
from.

\begin{aside}
So the question that predicts which half of the pipeline a clip change pays through is not ``grid
or objective'' but \textbf{does it distort sensitivity ratios, or merely scale them?} Scaling pays
in the encoder; distorting ratios pays in the allocator.
\end{aside}

This also corrects a generalisation made earlier in the campaign. The deep grid's $+0.09$
encoder-only result had been read as ``the encoder is not where the points are.'' That was right
about that change and wrong as a rule --- the second split was run anyway, out of routine rather
than expectation, and it returned $+0.70$. \textbf{Run the decomposition every time, especially
when a previous one makes it look like a formality.}

%% ============================================================ 8
\section{The mechanism behind the size advantage: the tie}
\label{sec:tie}

The clearest explanation of \emph{why} there are bytes to save is visible in GPTQ's own two
configurations --- same algorithm on both sides:

\begin{center}\small
\begin{tabular}{@{}lrrl@{}}
\toprule
GPTQ 3-bit & acc \% & bytes & note \\
\midrule
\code{ignore=["lm\_head"]} & 88.97 & 1{,}558{,}371{,}533 & tie left fp16 $=$ 65.3\pct{} of stored bits \\
\code{ignore=[]} & 88.03 & 741{,}475{,}927 & tie quantized uniformly at 3 bits \\
\bottomrule
\end{tabular}
\end{center}

The difference is \best{$+0.94$, 106/56, $p = 0.0001$, CI $[+0.47, +1.41]$}. \textbf{Quantizing
the tied embedding at 3 bits costs GPTQ a full point.}

Compare the 4-bit RTN pair measured in Phase 1: 0.13 points at $p = 0.48$, \emph{not separated},
for a 43\pct{} byte saving. \textbf{The convention is one fact at 4 bits and a different one at
3} --- at 4 bits the tie survives quantization and \code{ignore=["lm\_head"]} is an unexamined
default worth dropping for free; at 3 bits the tie is where the damage concentrates.

GPTQ has no per-row width mechanism, so its only two options are a 2.2\x{} larger checkpoint or
that full point. The row plan \code{8b:2048, 4b:8192, 2b:*} is a third option neither reaches, and
metadata is one fp16 scale and offset per row per group regardless of how the rows are split ---
so the granularity is nearly free. \textbf{That asymmetry is the size advantage.}

Note that this is the \emph{same} asymmetry the body exploits in \S\ref{sec:granularity}: the
format prices metadata per row per group either way, so a per-row width vector is close to free
wherever it is applied, and the only question is whether the ordering handed to it is any good.

%% ============================================================ 9
\section{Claim discipline: what is a win and what is a tie}
\label{sec:discipline}

\begin{center}\small
\begin{tabular}{@{}llr@{}}
\toprule
claim & numbers & $p$ \\
\midrule
\rowcolor{winbg}
\textbf{beats GPTQ at 4.5\pct{} fewer bytes} & \code{rb\_agg} 89.57 @ 708.1 MB vs 88.03 @ 741.5 MB & \best{$<0.0001$} \\
\rowcolor{winbg}
\textbf{beats GPTQ at 8.3\pct{} fewer bytes} & \code{rb\_680} 89.09 @ 680.0 MB vs 88.03 & \best{0.0026} \\
ties GPTQ at 16.4\pct{} fewer bytes & \code{rb\_620} 87.84 @ 620.0 MB vs 88.03, $-0.19$ & 0.64 \\
ties GPTQ's fp16-tie config at 45\pct{} of its size & 89.57 @ 708.1 MB vs 88.97 @ 1{,}558.4 MB, $+0.60$ & 0.069 \\
\bottomrule
\end{tabular}
\end{center}

\textbf{Only the first two rows survive as wins.}

The last row deserves its own paragraph, because it is exactly the situation in which a favourable
result invites over-claiming. \code{rb\_agg} 89.57\pct{} at 708{,}087{,}808 B against
\code{gptq\_3b} 88.97\pct{} at 1{,}558{,}371{,}533 B is \best{$+0.60$, 161/129, $p = 0.069$, CI
$[-0.03, +1.23]$}. That interval includes zero. The direction \emph{reversed} when the row body
landed --- it was $-0.43$ with the module body --- and it is tempting to promote on that basis. It
does not clear the threshold. The correct phrasing is \emph{indistinguishable from the best GPTQ
configuration available, at 45\pct{} of its size} --- never ``beats.''

%% ============================================================ 10
\section{What this does not overturn}
\label{sec:notoverturn}

Phase 2 was built and measured on \textbf{one panel}: Qwen3.5-2B-Base / CaseHOLD. The new recipe
has not been run on Mistral-7B / Banking77 or on Qwen / Banking77. The Phase 1 scorecard rows for
those panels stand as written, and so does the campaign's broader reading that DynQuant's advantage
tracks the \emph{baselines'} inconsistency rather than its own.

Three specific reasons to expect less elsewhere, all established above:

\begin{itemize}[leftmargin=1.4em]
\item \textbf{The tie lever cannot transfer to an untied model.} On Mistral the head is 3.71\pct{}
of the checkpoint, not 27.05\pct, so the row-partitioned embedding --- the largest single lever
here at $+2.90$ --- has almost nothing to work with. Expect the Mistral gap to close by much less.

\item \textbf{The allocator lever needs budget pressure.} At 4.25 bits the knapsack is barely
constrained and the map is close to a fixed architectural prior; Phase 1 measured the signal as
worth nothing there ($+0.19$, $p = 0.15$). \textbf{These gains are a 3-bit-tier phenomenon by
construction.}

\item \textbf{The row-body lever is only as good as the signal on that panel.} Its shuffled control
\emph{loses} 1.28 points, so it is not a structural improvement that travels --- it is an
amplifier, and on a panel where the sensitivity table is less informative it amplifies less, or
amplifies the wrong thing. \textbf{The control has to be re-run wherever the recipe is re-run.}
\end{itemize}

Two further limits carry over from Phase 1 unchanged and are not addressed here: the evaluation is
classification accuracy on a fixed answer set, with no generative, perplexity or long-context
measurement; and the baselines remain one library's implementation at one untuned configuration
(256 calibration rows, \code{group\_size} 128, \code{dampening\_frac} 0.01). A better-tuned GPTQ
would narrow these margins.

\subsection{The honest one-line update to the Phase 1 scorecard}

\begin{aside}
\textbf{``Beats GPTQ at 3 b''} moves from \emph{fails on all three panels} to \emph{reversed on
Qwen/CaseHOLD --- $+1.54$ at 4.5\pct{} fewer bytes ($p < 0.0001$), still ahead at 8.3\pct{} fewer,
level at 16.4\pct{} fewer; untested with the new recipe on the other two.}
\end{aside}

That is a stronger claim than the section originally carried, and it is worth being explicit about
how it was reached, because \textbf{the sequence matters more than the endpoint}. The intermediate
arm was $+0.51$ at $p = 0.16$ and was written up as a tie, over several days, while it was the best
result available. Had it been promoted to ``beats'' then, the actual result --- which came from a
different mechanism landing underneath it, not from re-reading the same number more favourably ---
would have had nothing left to report.

%% ============================================================ 11
\section{What ships, and what is still open}

\subsection{Promote into \code{dynquant-core}}

Four components, in the order their measured value justifies:

\begin{enumerate}[leftmargin=1.4em]
\item \textbf{Row-partitioned tied embedding} ($+2.90$, 37\x{} the body's return per byte). Applies
to every tied-embedding model and is nearly free in metadata.
\item \textbf{Gauss--Newton sensitivity allocator} ($+1.07$, $p = 0.0002$) replacing the
rank-product scorer. The cleanest substitution in the phase.
\item \textbf{Per-row body allocation} ($+1.04$, $p = 0.0001$) --- \emph{shipped together with its
shuffled control as a required regression test}, because without the signal it is a 1.28-point
regression.
\item \textbf{$\mathbb{E}[x^2]$-weighted clip objective with the deep grid} (pooled $p = 0.0098$
for the objective, $+0.58$ for the grid through re-pricing).
\end{enumerate}

The allocator has one integration requirement that cost a day to find and is worth stating: it must
be given the tie's cost via \code{--reserve-embed-bits}. A map allocated to a whole-model average
puts the embedding at a narrow width and spends the savings on the body; substituting a row plan
afterwards does not hand the body's share back, and the checkpoint lands about 24 MB over target
with nothing reporting why.

\subsection{Open}

\begin{itemize}[leftmargin=1.4em]
\item \textbf{Ablate the row signal itself.} The row sensitivities are currently a
max-of-ranks composite. Gradient-only, frequency-only and rank-product variants have not been
compared at row granularity, and \S\ref{sec:granularity} predicts the comparison is now
\emph{more} sensitive than it was at module granularity, not less.
\item \textbf{Instrument the row descent.} The 560 MB collapse has no diagnostic. Role floors are
per module; a per-row analogue does not exist.
\item \textbf{Re-run the recipe on the other two panels}, with the shuffled control alongside it in
both cases.
\item \textbf{Finer partitions.} Per-group and per-column allocation are the natural next steps and
the format prices their metadata identically --- but each one needs its own shuffled control before
any number from it is quoted.
\end{itemize}

%% ============================================================ 12
\section{Addendum (2026-08-15): the asymmetric control, as a replicate}
\label{sec:addendum}

The checkpoint behind every number in \S1--\S11 no longer exists --- the box it was measured on
has since been recycled --- so the control this report held itself pending could not be run on
the same weights. It was run as a \textbf{replicate}: the same driver
(\code{p2replicate.sh}), the same recipe (LoRA $r=32$, lr $10^{-4}$, 2 epochs, Qwen3.5-2B-Base /
CaseHOLD, $n = 5{,}314$), a fresh fine-tune, on torch 2.12.0+cu130 / transformers 5.10.1 in place
of the original panel's torch 2.13 / transformers 5.14. Every comparison below is therefore
arm-vs-arm \emph{within this replicate's own run}, not against \S1--\S11's numbers, which stand
as a record of what was measured then, under code corrected since (\S\ref{sec:addendum-bytes}).

\subsection{The four-arm table}

\begin{center}\small
\begin{tabular}{@{}lrrrr@{}}
\toprule
arm & acc \% & correct/total & bytes & avg bits \\
\midrule
\rowcolor{winbg}
\code{p2\_rb\_agg} (DynQuant) & \best{89.57} & 4760/5314 & 708{,}087{,}552 & 3.0102 \\
\code{gptq\_3b\_head\_asym} (control) & 88.45 & 4700/5314 & 741{,}486{,}130 & 3.1522 \\
\code{gptq\_3b\_head} (symmetric) & 87.67 & 4659/5314 & 735{,}981{,}792 & 3.1288 \\
\code{bf16} (ceiling) & 89.42 & 4752/5314 & 3{,}763{,}692{,}048 & 16 \\
\bottomrule
\end{tabular}
\end{center}

\begin{center}\small
\begin{tabular}{@{}lrrrrr@{}}
\toprule
comparison & $\Delta$acc & discordant (b/c) & exact $p$ & 95\pct{} CI & $\Delta$bytes \\
\midrule
\rowcolor{winbg}
\code{rb\_agg} vs asymmetric control & \best{$+1.13$} & 155/95 & \best{0.000178} & $[+0.55,+1.71]$ & \best{$-4.50\pct{}$} \\
\rowcolor{winbg}
\code{rb\_agg} vs symmetric arm & \best{$+1.90$} & 212/111 & \best{\e{2.0}{-8}} & $[+1.24,+2.56]$ & \best{$-3.79\pct{}$} \\
\code{rb\_agg} vs bf16 ceiling & $+0.15$ & 81/73 & 0.573 (ns) & $[-0.31,+0.61]$ & $-81.19\pct{}$ \\
\bottomrule
\end{tabular}
\end{center}

\textbf{The held claim is settled.} The asymmetric control this report said was missing now
exists, and DynQuant beats it: $+1.13$ points at 4.5\pct{} \emph{fewer} bytes, not matched bytes.
The 1.71 points of fp16 headroom the original abstract worried about only cost the claim 0.41 of
its 1.54 points --- the real control closed less than a third of the gap it had room to close.
DynQuant also statistically ties the bf16 ceiling at $5.32\times$ fewer bytes, a comparison
\S1--\S11 did not make.

\subsection{What the original panel's own number turns out to say}
\label{sec:addendum-bytes}

A separate audit (\code{byte-accounting-zero-point.md}) found that both baseline stages priced
group metadata as \code{meta\_bits = 16 + bits} \emph{unconditionally} --- scale plus a zero
point --- even for symmetric arms, which store no zero point. Every symmetric GPTQ/RTN arm this
project has published was over-charged by that difference, about 0.7\pct{} of a 3-bit checkpoint.
That audit could not correct \S1--\S11's own \code{gptq\_3b\_head} figure, because its checkpoint
was already gone, but it predicted the direction: \emph{``the replicate's GPTQ arm will come in
$\sim$0.7\pct{} cheaper than the destroyed panel's did.''}

\begin{center}\small
\begin{tabular}{@{}lrr@{}}
\toprule
 & bytes & $\Delta$ vs the original panel's 741{,}475{,}927 B \\
\midrule
this replicate's symmetric arm & 735{,}981{,}792 & \best{$-0.741\pct{}$} \\
this replicate's asymmetric arm & 741{,}486{,}130 & $+0.0014\pct{}$ \\
\bottomrule
\end{tabular}
\end{center}

The prediction lands almost exactly: $-0.741\pct{}$ against a forecast of $\sim\!-0.7\pct{}$.
It also resolves something \S1--\S11 left unexplained: \code{gptq\_3b\_head} was fitted and
labelled \emph{symmetric}, yet its byte figure sits within 10{,}203 bytes (0.0014\pct{}) of this
replicate's true \emph{asymmetric} figure, and 5{,}494{,}135 bytes (0.75\pct{}) away from this
replicate's true symmetric figure. The original number was a symmetric arm's weights, accounted
through code that priced every arm as though it stored a zero point --- a symmetric-fit arm
wearing an asymmetric-sized byte count. It also explains why \S1's abstract asked for the control
to be run ``at 741{,}475{,}927 bytes'': under that accounting, symmetric and asymmetric were
expected to cost the same, because the code never priced the difference between them. They do
not. This replicate's own two controls show a genuine 0.75\pct{} gap between them, which is
exactly the term the fix restored.

\subsection{What this does not change}

No number in \S1--\S11 is edited in place --- the checkpoint behind it cannot be re-measured, and
its table stands as a record of what was true under code since corrected. The four-lever
attribution in \S3 is untouched: every lever there was isolated at byte-identical budgets between
DynQuant arms only, and nothing about the GPTQ control bears on those deltas. Scope is still one
panel; \S\ref{sec:notoverturn} stands as written.

\vspace{1.5em}
\begin{center}
{\color{dqblue}\rule{0.5\textwidth}{0.8pt}}
\end{center}
\vspace{0.5em}

{\small\itshape
Every arm in this report was scored on the same 5{,}314 CaseHOLD items in the same order, every
$p$-value is an exact paired McNemar test on stored per-item hits, every byte-identical A/B is
byte-identical by construction rather than by tuning, and the winning arm was re-encoded and
audited for materialization before its number was believed. The two negative controls
(\code{rb\_shuf}, \code{tie\_u2\_ctl}) are reported at the same weight as the wins, because in this
phase they were the more informative measurements.\par}

\end{document}
